% ============================================================
\appendix
\section{The DM--CM Transformation Through the Lens of Linear Algebra}
\label{app:linalg}
% ============================================================

The transformation in \S2.2 can feel like a clever algebraic trick
specific to EMC.  It is not.  It is a direct application of
\emph{change of basis} — one of the central ideas in Gilbert Strang's
\textit{Introduction to Linear Algebra}.  This appendix unpacks the
four deepest connections.

% ---- A.1 ----
\subsection{The Same Circuit in New Coordinates — Strang Chapter~8}

Strang's core message on linear transformations is this: a linear map
is a fixed, coordinate-free object.  A \emph{matrix} is only its
\emph{representation} in a chosen basis.  Change the basis, and the
same map is represented by a different matrix — but nothing physical
has changed.

For a linear map $A \colon V \to V$ acting \emph{within a single
vector space}, changing basis with matrix $P$ gives the
\textbf{similarity transformation}:
\[
    A_{\mathrm{new}} = P^{-1} A P.
\]
The same $P$ appears on both sides because domain and codomain are
the same space — they are transformed by the same matrix.

\textbf{The circuit case is different — and this is the key point.}
The impedance matrix $\mathbf{Z}$ does not map a space to itself.
It maps \emph{currents} to \emph{voltages} — two physically distinct
spaces:
\[
    \mathbf{Z}
    \colon \underbrace{\mathbb{R}^N}_{\text{current space}}
    \;\longrightarrow\;
    \underbrace{\mathbb{R}^N}_{\text{voltage space}}.
\]
For a map $A \colon V \to W$ between \emph{two different spaces},
changing basis in the domain by $P$ and in the codomain by $Q$
\emph{independently} gives the correct formula:
\[
    \boxed{A_{\mathrm{new}} = Q\,A\,P^{-1}.}
\]
Here $P = \TiN$ (new basis for current space) and
$Q = \TvN$ (new basis for voltage space), so:
\[
    \underbrace{\mathbf{Z}}_{\text{per-line basis}}\,\mathbf{i}
    = \mathbf{v}_{nP}
    \quad\xrightarrow{\;P=\TiN,\;Q=\TvN\;}\quad
    \underbrace{\mathbf{Z}_{\mathrm{DCM}}}_{\text{DM--CM basis}}\,\iDCM
    = \vDCM,
    \qquad
    \mathbf{Z}_{\mathrm{DCM}} = \TvN\,\mathbf{Z}\,(\TiN)^{-1}.
\]
The similarity formula $P^{-1}AP$ is the special case $Q = P$ —
i.e.\ when domain and codomain are the same space transformed by the
\emph{same} matrix.  Here $Q \neq P$ because currents and voltages
obey different physical laws, forcing different basis-change matrices.
Why exactly $\TvN \neq \TiN$ is the subject of \S A.4.

\begin{statebox}[Strang's Framing, Applied Here]
``The whole point of a basis change is to find a coordinate system
in which the problem looks simpler.''\\[4pt]
Here, \emph{simpler} means: DM and CM appear as explicit, separate
rows in the equation — so the question ``does DM voltage drive CM
current?'' is answered by reading a single matrix entry, not by
solving $N$ equations and post-processing.
\end{statebox}

% ---- A.2 ----
\subsection{How Does Anyone Invent These Matrices? — The Basis Comes From the Problem}

The most common reaction to seeing $\TiN$ and $\TvN$ for the first
time is: \emph{how did the author know to use these?}

Strang answers the analogous question about eigenvectors by saying:
the right basis for a problem is not invented algebraically — it
emerges from the \emph{structure of the problem itself}.  For
eigenproblems, the structure is the matrix $A$.  For DM--CM
decomposition, the structure is \emph{what your instruments measure}.

\begin{physbox}[The Matrices Write Themselves From Measurement Rules]
\begin{itemize}
    \item A \textbf{clamp meter} looped around all $N$ wires reads
          $i_1 + i_2 + \cdots + i_N$.  That IS the bottom row of
          $\TiN$: $\;[\,1,\;1,\;\ldots,\;1\,]$.

    \item A \textbf{differential current probe} on lines $m$ and $n$
          reads $\tfrac{1}{2}(i_m - i_n)$.  That IS the corresponding
          DM row of $\TiN$: $\;[\,\ldots,\;\tfrac{1}{2},\;\ldots,\;
          {-\tfrac{1}{2}},\;\ldots\,]$.
\end{itemize}
Write down every measurement rule as a row.  Stack the rows.  You
have $\TiN$.  The matrix was never invented — it was \emph{read off
from the measurement setup}.
\end{physbox}

There is also a deeper algebraic reason the DM--CM directions are
natural.  Define the CM vector $\mathbf{c} = [1, 1, \ldots, 1]^T$.
\begin{itemize}
    \item The \textbf{CM subspace} is $\operatorname{span}\{\mathbf{c}\}$
          — all lines oscillate together, in the same direction.
    \item The \textbf{DM subspace} is the null space of $\mathbf{c}^T$
          — vectors whose entries sum to zero (pure differences, no
          common component).
\end{itemize}
By Strang's Fundamental Theorem of Linear Algebra (Chapter~3), the
row space and null space of any matrix are orthogonal complements.
Here: the CM subspace ($1$-dimensional) and the DM subspace
($(N{-}1)$-dimensional) are orthogonal complements in $\mathbb{R}^N$
and together span all of $\mathbb{R}^N$.  Every per-line current
vector splits \emph{uniquely} into a CM part and a DM part — this
uniqueness is what makes the transformation invertible (recall \S1.6).

% ---- A.3 ----
\subsection{Diagonalization, Coupling, and Eigenvectors — Strang Chapter~6}

Strang's diagonalization theorem: a matrix $A$ is diagonal in basis
$P$ — i.e.\ $P^{-1}AP = \Lambda$ — \emph{if and only if} the
columns of $P$ are eigenvectors of $A$.

Now apply this directly to the two cases of this paper.

\textbf{Symmetric circuit ($k = 1$).}\;
$\mathbf{S}_N = \tfrac{1}{C}\mathbf{I}_N$, so
$\mathbf{Z} = \bigl(R + pL + \tfrac{1}{pC}\bigr)\mathbf{I}_N
             = \lambda\,\mathbf{I}_N$.
A scalar multiple of the identity has \emph{every} vector as an
eigenvector.  In particular, the DM--CM basis vectors are
eigenvectors of $\mathbf{Z}$.  Therefore $\mathbf{Z}_{\mathrm{DCM}}$
is diagonal: the DM and CM equations are fully decoupled.

\textbf{Asymmetric circuit ($k \neq 1$).}\;
$\mathbf{S}_N = \operatorname{diag}(1/C, \ldots, 1/C, 1/(kC))$,
which is \emph{not} proportional to $\mathbf{I}_N$.  The DM--CM
basis vectors are \textbf{not} eigenvectors of $\mathbf{Z}$.
Therefore $\mathbf{Z}_{\mathrm{DCM}}$ is not diagonal — off-diagonal
entries appear.

\begin{eqnbox}[Strang's Lesson, Inverted]
Strang uses diagonalization to \emph{simplify} a matrix by finding
its eigenvector basis.  Here we do the opposite: we deliberately
\textbf{choose a non-eigenvector basis} (DM--CM) because that is what
EMC instruments measure.\\[4pt]
The off-diagonal entries of $\mathbf{Z}_{\mathrm{DCM}}$ are not a
failure to diagonalize — they are the \textbf{coupling we came to
find}, made explicitly readable by the change of coordinates.  A DM
voltage source produces CM current if and only if these off-diagonal
entries are non-zero.
\end{eqnbox}

This also tells you immediately \emph{where} the coupling lives: only
in the rows and columns that involve line $N$ (the asymmetric line),
because that is the only entry of $\mathbf{S}_N$ that differs from
the rest.  The algebra does not hide this — once you are in the right
coordinate system, you can see it directly.

% ---- A.4 ----
\subsection{Why $\TvN \neq \TiN$ — KVL and KCL Are Different Operations}

Section~A.1 established that the correct formula for a linear map
between two spaces is $A_{\mathrm{new}} = Q\,A\,P^{-1}$, not
$P^{-1}AP$.  The remaining question is: \emph{why} does this circuit
require $Q = \TvN \neq P = \TiN$?  The answer is physical.

Voltages and currents are governed by \emph{different physical laws}
that produce different natural operations:

\begin{itemize}
    \item \textbf{Currents obey KCL} (sum at a node).  CM current is
          therefore a \emph{sum}: $\icm = i_1 + \cdots + i_N$.
          DM current is a \emph{difference}: $i_{mn} =
          \tfrac{1}{2}(i_m - i_n)$.

    \item \textbf{Voltages obey KVL} (sum around a loop).  CM voltage
          is an \emph{average}: $\vcm = \tfrac{1}{N}\sum v_{nP}$.
          DM voltage is a \emph{full line-to-line difference}:
          $\vdm{m}{n} = v_{mP} - v_{nP}$.
\end{itemize}

Notice the asymmetry: DM current carries a factor of $\tfrac{1}{2}$
but DM voltage does not.  This is why $\TvN \neq \TiN$.  The factor
choices are fixed by two requirements: (1) the definitions must match
what instruments physically measure, and (2) the transformation must
be invertible so that $i_n$ and $v_{nP}$ can be recovered from
$\iDCM$ and $\vDCM$ (as shown in \S1.6).

Strang discusses adjoint operators and dual spaces (Chapter~6) and
shows that when you change coordinates for one quantity in a bilinear
pairing, the other quantity transforms by the \emph{transpose-inverse}
to preserve the pairing.  For $N = 2$ lines, this adjoint
relationship holds exactly: $\TvN[2] = (\TiN[2])^{-T}$, and the
DM--CM power identity $v_{12}\,I_{12} + \vcm\,\icm =
v_{1P}\,i_1 + v_{2P}\,i_2$ is satisfied.  For $N > 2$ the
definitions are fixed by measurement convention, and the clean
adjoint relationship does not extend — a subtlety the paper does not
need to resolve because it works entirely at the level of the
transformed impedance matrix $\mathbf{Z}_{\mathrm{DCM}}$.

\begin{statebox}[Four Strang Connections — Summary]
\centering
\begin{tabular}{@{}p{1cm}p{3cm}p{9cm}@{}}
    \toprule
    \textbf{\#} & \textbf{Strang concept} & \textbf{Role here} \\
    \midrule
    A.1 & Change of basis (Ch.~8) &
          $\mathbf{Z}_{\mathrm{DCM}} = \TvN\,\mathbf{Z}\,(\TiN)^{-1}$
          is same circuit, new coordinates \\
    \addlinespace
    A.2 & Right basis from problem structure (Ch.~3, 6) &
          Rows of $\TiN$, $\TvN$ are instrument measurement rules \\
    \addlinespace
    A.2 & Fundamental Theorem — null space $\perp$ row space (Ch.~3) &
          DM subspace $\perp$ CM subspace; decomposition is unique \\
    \addlinespace
    A.3 & Diagonalization (Ch.~6) &
          Symmetric $\Rightarrow$ diagonal $\mathbf{Z}_{\mathrm{DCM}}$;
          asymmetric $\Rightarrow$ off-diagonal = coupling \\
    \addlinespace
    A.4 & Adjoint / dual spaces (Ch.~6) &
          KVL $\neq$ KCL $\Rightarrow$ $\TvN \neq \TiN$ \\
    \bottomrule
\end{tabular}
\end{statebox}

